\documentclass{article}

\usepackage[outputdir=build]{Hypo-Note}

\title{Hypoxanthine-LaTeX Manual}
\author{Hypoxanthine-LaTeX Project}
\date{v0.4.0}

\begin{document}
\maketitle
\tableofcontents

\section{Overview}

Hypoxanthine-LaTeX is a modular, engineering-friendly LaTeX ecosystem.

\begin{itemize}
	\item \textbf{Classes} provide \emph{entrypoints} and user-facing options: \texttt{Hypo-Note}, \texttt{Hypo-Sheet}.
	\item \textbf{Core} provides atomic capabilities: Fonts, Colors, Base, Math.
	\item \textbf{Modules} provide optional plugins: Box, Refs, Code, \dots
\end{itemize}

The public interface is defined by \texttt{FEATURES.md}; this manual is a more structured explanation and set of examples.

\section{Quick Start}

In your document, load an entry package by \emph{package name}:

\begin{verbatim}
\documentclass{article}
\usepackage[outputdir=build]{Hypo-Note}

\begin{document}
Hello.
\end{document}
\end{verbatim}

If you build via this repository's Makefile template, it injects \texttt{TEXINPUTS} so that \texttt{\textbackslash usepackage\{Hypo-Note\}} resolves to the local \texttt{sty/} directory.

\section{Entrypoints (Classes)}

\subsection{Hypo-Note}

	exttt{Hypo-Note} is the note entry (comfortable reading). It loads core capabilities and (optionally) modules.

\subsection{Hypo-Sheet}

	exttt{Hypo-Sheet} is the sheet/cheatsheet entry (compact layout). In the current stage it shares the same public options and loads the same core/modules.

\section{Global Options}

Both entrypoints accept these options:

\begin{itemize}
	\item \texttt{indent=true/false}: whether paragraphs are indented.
	\item \texttt{shorthand=true/false}: shorthand commands (enabled by default).
	\item \texttt{boxes=true/false}: whether box environments are loaded (enabled by default).
	\item \texttt{refs=true/false}: whether the references module (hyperref + cleveref) is loaded (enabled by default).
	\item \texttt{outputdir=\dots}: exported by Base as \texttt{\textbackslash FinalOutputDir}, and used by modules that need an output directory.
\end{itemize}

Example:

\begin{verbatim}
\usepackage[shorthand=false,boxes=false,refs=false,outputdir=build]{Hypo-Note}
\end{verbatim}

\section{Core}

\subsection{Fonts}

Fonts are configured for XeLaTeX and CJK via \texttt{xeCJK}.

\subsection{Colors}

Project colors are defined in \texttt{Hypo-Colors}. You can use them directly:

\begin{verbatim}
{\color{HypoBlueDark}Text}
\end{verbatim}

\subsection{Base (Output Directory)}

Base exports the final output directory:\quad \texttt{\FinalOutputDir = \FinalOutputDir}.

\subsection{Math and Shorthand}

Shorthand is enabled by default. Common commands include:

\begin{itemize}
	\item Text: \texttt{\textbackslash TX\{\dots\}}, \texttt{\textbackslash TBF\{\dots\}}
	\item Math: \texttt{\textbackslash MB\{\dots\}}, \texttt{\textbackslash MC\{\dots\}}, \texttt{\textbackslash BS\{\dots\}}
	\item Delimiters: \texttt{\textbackslash Abs\{\dots\}}
\end{itemize}

Demo:
\[
	\MB{x} + \MC{F}(\BS{\theta})\quad\text{and}\quad \Abs{-3}=3.
\]

\section{Boxes (Module: Hypo-Box)}

When \texttt{boxes=true} (default), the following environments are available:
	exttt{definition}, \texttt{example}, \texttt{note}.

\subsection{Title and Explicit Label}

All box environments take two mandatory arguments:

\begin{verbatim}
\begin{definition}{Title}{label}
	...
\end{definition}
\end{verbatim}

The actual label is automatically prefixed:

\begin{itemize}
	\item definition: \texttt{def:}\quad (e.g. \texttt{def:group})
	\item example: \texttt{ex:}\quad (e.g. \texttt{ex:integers})
	\item note: \texttt{note:}\quad (e.g. \texttt{note:tips})
\end{itemize}

If you do not need references, pass an empty label argument \texttt{\{\}}.

\subsection{Numbering Strategy (Advanced)}

The box module supports a \texttt{numbering=none/global/section/chapter} option (default: \texttt{section}).

If you want to control this option, disable boxes in the entrypoint and load \texttt{Hypo-Box} manually:

\begin{verbatim}
\usepackage[boxes=false]{Hypo-Note}
\usepackage[numbering=global]{Hypo-Box}
\end{verbatim}

\subsection{Style Overrides}

You can override styles via \texttt{tcolorbox} hooks:

\begin{verbatim}
	cbset{hypo definition box/.style={colframe=HypoRedDark,colback=HypoRedLight}}
\end{verbatim}

\subsection{Live Examples}

\begin{definition}{Group}{group}
	A group is a set with an operation satisfying closure, associativity, identity and inverse.
\end{definition}

\begin{example}{Integers}{integers}
	$(\mathbb{Z}, +)$ is a group.
\end{example}

\begin{note}{Tip}{}
	If you don\textquotesingle t need references, leave the label empty.
\end{note}

\section{References (Module: Hypo-Refs)}

When \texttt{refs=true} (default), \texttt{hyperref} and \texttt{cleveref} are loaded.
Use \texttt{\textbackslash cref\{\dots\}} to produce formatted references, for example:

\begin{verbatim}
See \cref{def:group}.
\end{verbatim}

In this project the default output is in the form "Definition: 0.1", "Example: 0.1", "Note: 0.1".

\section{Planned Features}

Some interfaces are reserved in \texttt{FEATURES.md} but not implemented yet, for example:\quad \texttt{\textbackslash img\{path\}\{caption\}}.

\end{document}
